% Options for packages loaded elsewhere
\PassOptionsToPackage{unicode}{hyperref}
\PassOptionsToPackage{hyphens}{url}
\documentclass[
  a4paper,
]{article}
\usepackage{xcolor}
\usepackage[margin=25mm]{geometry}
\usepackage{amsmath,amssymb}
\setcounter{secnumdepth}{-\maxdimen} % remove section numbering
\usepackage{iftex}
\ifPDFTeX
  \usepackage[T1]{fontenc}
  \usepackage[utf8]{inputenc}
  \usepackage{textcomp} % provide euro and other symbols
\else % if luatex or xetex
  \usepackage{unicode-math} % this also loads fontspec
  \defaultfontfeatures{Scale=MatchLowercase}
  \defaultfontfeatures[\rmfamily]{Ligatures=TeX,Scale=1}
\fi
\usepackage{lmodern}
\ifPDFTeX\else
  % xetex/luatex font selection
\fi
% Use upquote if available, for straight quotes in verbatim environments
\IfFileExists{upquote.sty}{\usepackage{upquote}}{}
\IfFileExists{microtype.sty}{% use microtype if available
  \usepackage[]{microtype}
  \UseMicrotypeSet[protrusion]{basicmath} % disable protrusion for tt fonts
}{}
\makeatletter
\@ifundefined{KOMAClassName}{% if non-KOMA class
  \IfFileExists{parskip.sty}{%
    \usepackage{parskip}
  }{% else
    \setlength{\parindent}{0pt}
    \setlength{\parskip}{6pt plus 2pt minus 1pt}}
}{% if KOMA class
  \KOMAoptions{parskip=half}}
\makeatother
\usepackage{longtable,booktabs,array}
\newcounter{none} % for unnumbered tables
\usepackage{calc} % for calculating minipage widths
% Correct order of tables after \paragraph or \subparagraph
\usepackage{etoolbox}
\makeatletter
\patchcmd\longtable{\par}{\if@noskipsec\mbox{}\fi\par}{}{}
\makeatother
% Allow footnotes in longtable head/foot
\IfFileExists{footnotehyper.sty}{\usepackage{footnotehyper}}{\usepackage{footnote}}
\makesavenoteenv{longtable}
\usepackage{graphicx}
\makeatletter
\newsavebox\pandoc@box
\newcommand*\pandocbounded[1]{% scales image to fit in text height/width
  \sbox\pandoc@box{#1}%
  \Gscale@div\@tempa{\textheight}{\dimexpr\ht\pandoc@box+\dp\pandoc@box\relax}%
  \Gscale@div\@tempb{\linewidth}{\wd\pandoc@box}%
  \ifdim\@tempb\p@<\@tempa\p@\let\@tempa\@tempb\fi% select the smaller of both
  \ifdim\@tempa\p@<\p@\scalebox{\@tempa}{\usebox\pandoc@box}%
  \else\usebox{\pandoc@box}%
  \fi%
}
% Set default figure placement to htbp
\def\fps@figure{htbp}
\makeatother
\setlength{\emergencystretch}{3em} % prevent overfull lines
\providecommand{\tightlist}{%
  \setlength{\itemsep}{0pt}\setlength{\parskip}{0pt}}
\usepackage{bookmark}
\IfFileExists{xurl.sty}{\usepackage{xurl}}{} % add URL line breaks if available
\urlstyle{same}
\hypersetup{
  hidelinks,
  pdfcreator={LaTeX via pandoc}}

\author{}
\date{}

\begin{document}

\section{Geometric Rapid Expansion Is Specific to Unstable
Self-Consistent Closures: Causal Discrimination of Onset Modes via
Generic Zero-Closure Initial
States}\label{geometric-rapid-expansion-is-specific-to-unstable-self-consistent-closures-causal-discrimination-of-onset-modes-via-generic-zero-closure-initial-states}

\textbf{Author:} Noriaki Kihara (WF System Co., Ltd.) \textbf{Date:}
2026-08-05 \textbf{Version DOI:}
\href{https://doi.org/10.5281/zenodo.21798855}{10.5281/zenodo.21798855}
\textbf{Concept DOI:}
\href{https://doi.org/10.5281/zenodo.21798854}{10.5281/zenodo.21798854}

\begin{center}\rule{0.5\linewidth}{0.5pt}\end{center}

\subsection{Claims}\label{claims}

Paper 8 {[}3{]} causally removed the two explicit seeds from the time
evolution discovered in Paper 7 {[}2{]} --- long latency, geometric
rapid expansion, and spontaneous formation of a metastable
three-directional structure --- and established that the phenomenon does
not depend on the seeds. At the same time, Paper 8 left the causality of
the initial state itself unverified, as a stated boundary condition. The
experiments had consistently started from the self-consistent circularly
polarized eigenmode, i.e., a maximally coherent state in which all
relational waves ride on a single eigen-oscillation, and Paper 8
explicitly stated that ``a systematic control from generic states
satisfying only the closure condition is the next task'' (Paper 8 §9.6).
The present paper is the execution report of that control.

First, we built a generator that constructs generic zero-closure initial
states without contrivance. The raw material is Gaussian white noise,
and closure is not imposed afterwards by projection or correction; it
holds identically, term by term, through the construction of two
orthogonal equal-norm components. The removal of the DC component
(\(k=0\)) and, for even \(N\), the Nyquist component (\(k=N/2\)) is not
an additional assumption. The zero square sum of a single mode \(k\) is
proportional to \(\sum_n e^{4\pi ikn/N}\), so a mode with
\(2k\equiv 0 \pmod N\) can never close on its own. The first axiom
already denies these components the qualification to exist; the removal
is the enforcement of a theorem (§2).

Second, with the classification criteria fixed before execution, we
measured a dichotomy of onset modes (§3--§4). The series starting from
self-consistent closures are of the latency-burst type: after a latency
of 288--695 steps, \(\log f\) grows linearly over about 18 decades (N=5:
0.0494/step, coefficient of determination \(R^2=0.9999998\); N=40:
0.0350/step, \(R^2=0.998\)). The generic zero-closure white harmonic
sectors --- all 42 of them across N=5 and N=40 --- are of the immediate
type: zero floor residence, crossing at 11--16 steps, no amplification
epoch whatsoever, and an immediate transition to the metastable
three-directional structure (rank 4). There is not a single exception.

Third --- this is the central claim of the paper. Geometric
amplification is not a universal property of arbitrary closed states. It
is a phenomenon specific to the instability of the self-consistent
closure --- the single coherent condensate --- and no amplification
arises from a sea that already carries harmonic structure. The
metastable three-directional structure is a terminal point common to
both onset modes; the only fossil that distinguishes the onset mode is
the trace of amplification. Inside this model, therefore, a time
evolution bearing the trace of 18-decade exponential amplification and
an onset from white chaos are incompatible. The onset mode is neither an
assumption nor an interpretation; it is a proposition dynamically
decidable by the presence or absence of amplification. This proposition
is sharpened by the stability measurements of the sixth item into
``specific to unstable self-consistent closures'' --- the title adopts
that final form.

Fourth, whether amplification occurs is determined not by the provenance
of the state but by its structure alone (§6). The parent state obtained
by solving the self-consistency iteration with white noise as its seed
(N=40) shows the same latency-burst type as the deterministically
constructed control, with matching growth rates (0.0347 vs.~0.0350,
under 1\% difference). The parent of the same generator at N=5 converged
to a different eigen-solution and did not amplify out to 12000 steps.
What is established at this stage is therefore the necessity side: if
amplification occurs, the initial state is a self-consistent relative
equilibrium. This asymmetry is resolved by the stability measurements of
the sixth item, which establish that, within the tested range,
\textbf{an unstable relative equilibrium is the necessary and sufficient
condition}. What remains is a theoretical classification predicting,
from the structure of a solution, which solutions are stable and which
unstable (§8).

Fifth, the dichotomy admits a natural mechanistic reading (§5). The
current update rule carries no explicit generation vertex built from
amplitude products --- a fact mechanically verifiable from the structure
of the code. The self-consistent closure is a relative equilibrium of
the dynamics; geometric amplification is its unstable departure, and the
immediate departure of generic states is the ordinary motion of
non-equilibrium states. We state this reading explicitly as the linear
no-generation hypothesis and do not claim a proof of a general
no-generation theorem. That the white sea behaves as an already-unfolded
state from the start, and that only the condensate has room to amplify,
are two faces of this reading.

Sixth, we converted this mechanistic reading into measurements through
three verification experiments (§5.1--5.3). (i) The one-step residual is
\(\le2.5\times10^{-15}\) for the 4 parent states and
\(\ge1.8\times10^{-2}\) for the 82 non-equilibrium states (42 sectors +
40 generic states) --- a 13-order-of-magnitude separation with nothing
in between. ``Self-consistent closure = relative equilibrium'' is a
direct observable. (ii) The maximal growth rate of the tangent map at
the relative equilibria matches the prediction
\(\lambda_{\max}=\text{rate}_f/2\) exactly at N=5 (0.02468 vs.~0.02468)
and within 4\% at N=40, and the inert white-origin parent (N=5) turned
out to be a stable equilibrium with not a single unstable eigenvalue.
Hence, across all 86 tested states, the equivalence \textbf{latency
exponential amplification \(\iff\) unstable relative equilibrium} holds
without exception. (iii) In the sweep of the perturbation amplitude
\(\varepsilon\), \(t_{\mathrm{cross}}\) is linear in \(-\ln\varepsilon\)
(N=5: measured/predicted slope ratio 0.996) and saturates at the
natural-trajectory value for \(\varepsilon\lesssim3\times10^{-13}\).
This saturation dynamically measures the location of the intrinsic seed
of the seedless trajectory; its value
\(\varepsilon_{\mathrm{eff}}\approx3.4\times10^{-13}\) (N=5) is of the
same order as the parent's eigenmode residual \(2.14\times10^{-13}\) ---
a partial answer to Paper 8's open problem (i), the physical origin of
the first minute difference. Furthermore, all 40 generic zero-closure
states (direct random samples from the manifold and sector mixtures) are
of the immediate type, establishing that the dichotomy is not a
peculiarity of single sectors.

Seventh, we record the full process of refutation and correction (§7).
The first version of the generator satisfied closure through the
cancellation of cross pairs of harmonics \(k\) and \(N-k\), leaving the
self-paired sectors unclosed --- a defect. The first closure measurement
on the experimental side also used a wrong quadratic form. Both were
corrected through control experiments (bit reproduction of the original,
correction, re-measurement). The central dichotomy was observed in both
the defective first version and the corrected version, surviving the
bookkeeping correction. The result is not an artifact of a particular
implementation.

Finally, the experiments in this paper were not designed on the basis of
prior work. The starting point was an internal question about the
naturalness of the initial spectrum; a control was designed
independently, a result contradicting the working hypothesis was
obtained, and only then was the prior literature re-surveyed ---
revealing that the standard inflationary initial-condition problem is
exactly the point of controversy that this result touches (§10). The
process by which a self-consistent closure spontaneously forms out of a
white sea (the capture problem) is not observed within the scope of this
paper, and the generation of structure from the condensate lies outside
the linear dynamics, being treated in the series on inelastic vertices
derived from closure conservation {[}5{]}. The unresolved part of the
onset-mode problem is thereby confined to a single location: the capture
from chaos to condensate.

\begin{center}\rule{0.5\linewidth}{0.5pt}\end{center}

\subsection{1. System (Inherited
Definitions)}\label{system-inherited-definitions}

\textbf{Axiom 1 (nontrivial zero-square-sum closure)}
\(\sum_e Z_e^2=0,\ Z\neq0\). \textbf{Axiom 2 (finite recurrence)}
\(U^n=I\). Both axioms follow the foundational axiom system {[}1{]}.

The dynamical system is identical to Papers 7 {[}2{]} and 8 {[}3{]}. A
complex amplitude \(Z_e\) is placed on each edge \(e\) of the complete
graph \(K_N\) (\(M=N(N-1)/2\)); the state closes on \(Z\in\mathbb C^M\),
\(\lVert Z\rVert=1\), \(Z^{\mathsf T}Z=0\). The generator \(K\) is a
real antisymmetric matrix depending only on edge phases, and the update
is a Cayley transform after spectral-norm normalization
(\(\gamma=\tan(\pi/144)\)), preserving the norm exactly. The dynamics
code is the Paper 7 original, imported read-only with SHA-256 fixation
and never modified (§11).

The observable is the departure fraction from each series' \textbf{own
initial plane}. From the initial state \(v_0\) we form
\(p=\mathrm{Re}\,v_0/\lVert\mathrm{Re}\,v_0\rVert\) and \(q\)
(\(\mathrm{Im}\,v_0\) orthogonalized against \(p\)), and define

\[f(t)=\frac{\lVert Z-\Pi_{\mathrm{span}\{p,q\}}Z\rVert^2}{\lVert Z\rVert^2}\]

By construction \(f(0)=0\) for every series, so onset modes are compared
in the same coordinate. Crossing is the first step at which \(f>0.05\),
the pre-existing threshold of Paper 7.

\subsection{2. Construction of Generic Zero-Closure Initial
States}\label{construction-of-generic-zero-closure-initial-states}

\subsubsection{2.1 The generator (white, zero-closure, post-hoc harmonic
readout)}\label{the-generator-white-zero-closure-post-hoc-harmonic-readout}

The only theoretical input is \(N\). Each relational wave \(w_m\) is an
\(N\)-sample complex waveform built from the two orthonormal columns
\((q_1,q_2)\) obtained by QR decomposition of \(N\times2\) Gaussian
white noise:

\[w_m=\frac{q_1+iq_2}{\sqrt2}\]

The zero square sum of each row holds identically:

\[\sum_n w_m[n]^2=\lVert q_1\rVert^2-\lVert q_2\rVert^2+2i\,q_1\!\cdot\!q_2=0\]

Closure is not manufactured by corrections; it emerges structurally from
two orthogonal equal-norm components. Harmonics are not placed by the
generator; they are read out after saving via an N-point DFT.

\subsubsection{2.2 DC/Nyquist removal is a theorem, not an
assumption}\label{dcnyquist-removal-is-a-theorem-not-an-assumption}

The zero square sum of a single mode \(k\) is proportional to
\(\sum_n e^{4\pi ikn/N}\): it vanishes if \(2k\not\equiv0\pmod N\) and
does not vanish if \(2k\equiv0\pmod N\). That is, DC (\(k=0\)) and, for
even \(N\), Nyquist (\(k=N/2\)) can never close alone for any \(N\), and
under the first axiom they lack the qualification to exist. The
corrected generator (v3) projects the DC direction (the all-ones vector)
and, for even \(N\), the Nyquist direction (the alternating \(\pm1\)
vector) out of each noise column before QR. Since the projection
preserves Gaussianity, this is a rigorous implementation of ``white
noise on the admissible modes.'' The strengthened version (v4) further
makes the closure debts between \(|k|\) pairs (pair-product debt)
identically zero, term by term, through direct spectral construction on
admissible signed modes. The dynamical experiments of this paper use v3;
v4 was used to audit the classification of initial states.

\subsubsection{2.3 Injected series}\label{injected-series}

Three kinds of initial state are injected into the same dynamics with
the same measurement.

\begin{enumerate}
\def\labelenumi{\arabic{enumi}.}
\tightlist
\item
  \textbf{control}: the self-consistent circularly polarized eigenmode
  returned by \texttt{make\_parent} of Papers 7--8 (deterministic
  construction).
\item
  \textbf{white-origin parent}: the eigenmode obtained by solving the
  self-consistency iteration with white noise as its initial value (the
  v3 generator's \texttt{parent\_vector}). Its structure ---
  self-consistent closure --- is the same as the control; only its
  provenance (random-noise origin) differs.
\item
  \textbf{white harmonic sectors}: states \(Z_k\) obtained by extracting
  only the harmonic-\(k\) content of the white sea as \(M\) relational
  amplitudes (all \(k\) with \(2k\not\equiv0\pmod N\); 4 for N=5, 38 for
  N=40, 42 in total). The field closure of each sector is exactly zero
  (measured: 0.0 for all 42). These are the ``generic states satisfying
  only the closure condition, with no eigenmode condition imposed.''
\end{enumerate}

\subsection{3. Classification Criteria (Fixed Before
Execution)}\label{classification-criteria-fixed-before-execution}

The following were fixed before execution, with post-hoc changes
prohibited.

\begin{itemize}
\tightlist
\item
  \textbf{Latency} \(t_{\mathrm{launch}}=\max\{t: f(t)<10^{-20}\}\) (the
  last time on the floor; 0 if \(f\ge10^{-20}\) from the start).
\item
  \textbf{Geometric criterion}: least-squares fit of \(\ln f\) vs.~\(t\)
  in the window \(f\in[10^{-20},10^{-2}]\); the window must span at
  least 6 decades with \(R^2>0.99\).
\item
  \textbf{Immediate criterion}: \(f(2)>10^{-6}\) (finite departure
  within 2 steps).
\item
  \textbf{Classes}: latency-burst type = \(t_{\mathrm{launch}}>100\) and
  geometric / immediate type = \(f(2)>10^{-6}\) / anything satisfying
  neither is recorded as ``other'' together with its raw data.
\end{itemize}

The measurement records \(f(t)\) at every step (\(T=6000\)). The only
random numbers are the generator seeds (N=5: 2, N=40: 1) and the
auxiliary injection seeds (parent 91000, sector \(92000+k\)), all
recorded.

\subsection{4. Result: the Dichotomy of Onset
Modes}\label{result-the-dichotomy-of-onset-modes}

\emph{Table 1: Classification of all series (canonical sources:
\texttt{paper8\_em9r\_profile\_N00005\_v1.json} /
\texttt{N00040\_v1.json})}

{\def\LTcaptype{none} % do not increment counter
\begin{longtable}[]{@{}
  >{\raggedright\arraybackslash}p{(\linewidth - 14\tabcolsep) * \real{0.1250}}
  >{\raggedright\arraybackslash}p{(\linewidth - 14\tabcolsep) * \real{0.1250}}
  >{\raggedleft\arraybackslash}p{(\linewidth - 14\tabcolsep) * \real{0.1250}}
  >{\raggedleft\arraybackslash}p{(\linewidth - 14\tabcolsep) * \real{0.1250}}
  >{\raggedleft\arraybackslash}p{(\linewidth - 14\tabcolsep) * \real{0.1250}}
  >{\raggedleft\arraybackslash}p{(\linewidth - 14\tabcolsep) * \real{0.1250}}
  >{\raggedleft\arraybackslash}p{(\linewidth - 14\tabcolsep) * \real{0.1250}}
  >{\raggedleft\arraybackslash}p{(\linewidth - 14\tabcolsep) * \real{0.1250}}@{}}
\toprule\noalign{}
\begin{minipage}[b]{\linewidth}\raggedright
Series
\end{minipage} & \begin{minipage}[b]{\linewidth}\raggedright
Class
\end{minipage} & \begin{minipage}[b]{\linewidth}\raggedleft
Latency
\end{minipage} & \begin{minipage}[b]{\linewidth}\raggedleft
Crossing
\end{minipage} & \begin{minipage}[b]{\linewidth}\raggedleft
Decades
\end{minipage} & \begin{minipage}[b]{\linewidth}\raggedleft
Rate /step
\end{minipage} & \begin{minipage}[b]{\linewidth}\raggedleft
\(R^2\)
\end{minipage} & \begin{minipage}[b]{\linewidth}\raggedleft
\(f(2)\)
\end{minipage} \\
\midrule\noalign{}
\endhead
\bottomrule\noalign{}
\endlastfoot
control N=5 & latency burst & 288 & 1166 & 17.97 & 0.0494 & 0.9999998 &
\(2.5\times10^{-29}\) \\
control N=40 & latency burst & 468 & 2011 & 17.99 & 0.0350 & 0.9981 &
\(2.0\times10^{-30}\) \\
white-origin parent N=40 & latency burst & 695 & 2167 & 17.98 & 0.0347 &
0.9997 & \(5.2\times10^{-32}\) \\
white-origin parent N=5 & inert (other) & --- (stays on floor) & none
(\(t\le12000\)) & 0 & --- & --- & \(3.0\times10^{-32}\) \\
white sectors N=5 (4) & immediate 4/4 & 0 & 11--16 & --- & --- & --- &
\(0.95\)--\(1.6\times10^{-3}\) \\
white sectors N=40 (38) & immediate 38/38 & 0 & 13--14 & --- & --- & ---
& \(1.25\)--\(1.46\times10^{-3}\) \\
\end{longtable}
}

\begin{figure}
\centering
\pandocbounded{\includegraphics[keepaspectratio,alt={N=5 rise profiles}]{figs/fig_em9r_profile_N00005.png}}
\caption{N=5 rise profiles}
\end{figure}

\emph{Figure 1: Rise profiles at N=5 (semilogy). Gray = control, black =
white-origin parent, blue = the 4 white sectors. Only the control shows
latency followed by 18 decades of linear growth (geometric
amplification). The sectors depart immediately with zero latency.}

\begin{figure}
\centering
\pandocbounded{\includegraphics[keepaspectratio,alt={N=40 rise profiles}]{figs/fig_em9r_profile_N00040.png}}
\caption{N=40 rise profiles}
\end{figure}

\emph{Figure 2: Rise profiles at N=40. The white-origin parent (black)
shows a latency burst of the same form as the control (gray), with
growth rates matching to under 1\%. All 38 sectors (blue) are
immediate.}

Three points stand out.

\begin{enumerate}
\def\labelenumi{\arabic{enumi}.}
\tightlist
\item
  \textbf{The presence or absence of an amplification epoch
  dichotomizes.} The three latency-burst series ride a single linear law
  of \(\log f\) over about 18 decades (best \(R^2=0.9999998\)). The 42
  immediate series contain no interval that could be called exponential
  growth: zero floor residence, finite departure to \(f\sim10^{-3}\)
  within 2 steps, crossing at 11--16 steps. Not a single intermediate
  type appeared.
\item
  \textbf{The terminal point is shared.} All 42 white sectors transition
  after crossing to the metastable three-directional structure (rank 4;
  late-time \(f\): 0.73--0.83 for N=5, 0.995--0.999 for N=40; canonical
  source: \texttt{paper8\_em9r\_v3\_result\_*.json}). The
  three-directional structure of Papers 7--8 is not specific to coherent
  onsets. What is specific is \textbf{the road there} --- the trace of
  amplification.
\item
  \textbf{The burst growth rate does not depend on preparation.} The
  N=40 growth rates of the control (deterministic construction) and the
  white-origin parent (random origin) agree at 0.0350 vs.~0.0347. The
  amplification rate is a property of the dynamics and the equilibrium
  structure, not of the state's provenance.
\end{enumerate}

\subsection{5. Mechanism: Amplification as Instability of a Relative
Equilibrium}\label{mechanism-amplification-as-instability-of-a-relative-equilibrium}

\textbf{Absence of an explicit generation vertex (a mechanically
verifiable fact).} One step of the current engine is a Cayley-type
unitary transform by the generator \(K(\theta)\) frozen at the current
phase configuration \(\theta=\arg Z\), and carries no explicit
generation vertex built from amplitude products of \(Z\). This is
verifiable from the structure of the code.

\textbf{Linear no-generation hypothesis (interpretation).} Each step is
linear with \(\theta\) frozen, but the full map \(F(Z)=U(\arg Z)\,Z\) is
state-dependent and hence nonlinear, so that iteration creates no new
content does not follow from the above fact alone. This paper interprets
the observed geometric amplification not as generation of new content
but as unstable departure from a relative equilibrium. This is a working
hypothesis explaining the dichotomy of this paper; no proof of a general
no-generation theorem is claimed. A proof of an invariant subspace, or a
direct derivation of the update equations of the harmonic coefficients,
is the task of the next stage.

Under this hypothesis the dichotomy reads as follows.

\begin{itemize}
\tightlist
\item
  The self-consistent closure is a \textbf{relative equilibrium} of the
  dynamics (fixed up to a uniform phase rotation). On the equilibrium,
  \(f\) does not strictly grow; departure occurs only through the linear
  instability of the equilibrium growing exponentially from the
  numerical floor (eigenmode residual \(\sim10^{-13}\) plus rounding).
  This is the identity of the latency and of the ensuing 18 decades of
  linear \(\log f\).
\item
  Generic zero-closure states are not equilibria. One step of the map
  immediately imparts a finite-angle rotation, so there is neither floor
  residence nor an exponential interval, and the state leaves its
  initial plane in \(O(10)\) steps. \textbf{The amplification epoch is
  absent because the ``small departure from equilibrium'' that would
  have to grow does not exist in the first place.}
\end{itemize}

Thus the measurement ``geometric amplification is specific to the
coherent condensate'' is not an accidental contrast; it is naturally
explained by the difference between relative equilibria and
non-equilibrium initial states. From the same viewpoint, the series of
negative results up to Paper 8 (no internal harmonic generation, no
precursor of statistical distinction) is consistent with the absence of
an explicit generation vertex --- although deriving them rigorously as
corollaries requires proving the hypothesis above.

This reading generates three verifiable predictions. (1) The one-step
residual \(r(Z_0)=\min_\phi\lVert F(Z_0)-e^{i\phi}Z_0\rVert\) is
numerically zero for parent states, finite for non-equilibrium states,
with nothing in between. (2) The logarithm of the largest tangent-map
eigenvalue \(\lambda_{\max}\) at a relative equilibrium equals half the
measured burst rate, \(\text{rate}_f/2\), since \(f\) is an
amplitude-squared quantity; the inert parent has no unstable eigenvalue.
(3) Injecting a perturbation amplitude \(\varepsilon\), the growth time
from \(f(0)\approx\varepsilon^2\) to \(f_{\mathrm{cross}}\) obeys the
logarithmic law
\(t_{\mathrm{cross}}\approx\text{const}+(-2\ln\varepsilon)/\text{rate}_f\).
The following three experiments (ONS-1/2/3, criteria and predictions
fixed before execution at the head of each script) tested these.

\subsubsection{5.1 Verification 1: One-Step Residual and Tangent-Map
Spectrum
(ONS-2)}\label{verification-1-one-step-residual-and-tangent-map-spectrum-ons-2}

\emph{Table 2: One-step residuals and tangent spectra (canonical source:
\texttt{onset\_equilibrium\_residual\_N*.json})}

{\def\LTcaptype{none} % do not increment counter
\begin{longtable}[]{@{}
  >{\raggedright\arraybackslash}p{(\linewidth - 8\tabcolsep) * \real{0.2000}}
  >{\raggedleft\arraybackslash}p{(\linewidth - 8\tabcolsep) * \real{0.2000}}
  >{\raggedleft\arraybackslash}p{(\linewidth - 8\tabcolsep) * \real{0.2000}}
  >{\raggedleft\arraybackslash}p{(\linewidth - 8\tabcolsep) * \real{0.2000}}
  >{\raggedleft\arraybackslash}p{(\linewidth - 8\tabcolsep) * \real{0.2000}}@{}}
\toprule\noalign{}
\begin{minipage}[b]{\linewidth}\raggedright
State
\end{minipage} & \begin{minipage}[b]{\linewidth}\raggedleft
Residual \(r\)
\end{minipage} & \begin{minipage}[b]{\linewidth}\raggedleft
\(\lambda_{\max}\) measured
\end{minipage} & \begin{minipage}[b]{\linewidth}\raggedleft
\(\lambda_{\max}\) predicted
\end{minipage} & \begin{minipage}[b]{\linewidth}\raggedleft
Unstable eigenvalues
\end{minipage} \\
\midrule\noalign{}
\endhead
\bottomrule\noalign{}
\endlastfoot
control parent N=5 & \(2.5\times10^{-15}\) & \textbf{0.02468} & 0.02468
& 4 \\
control parent N=40 & \(7.1\times10^{-16}\) & 0.01824 & 0.01750 & 383 \\
white-origin parent N=40 & \(1.0\times10^{-16}\) & 0.01815 & 0.01734 &
389 \\
white-origin parent N=5 (inert) & \(1.6\times10^{-16}\) &
\textbf{0.00000} & \(\le10^{-3}\) & \textbf{0} \\
white sectors N=5 (4) & \(2.1\)--\(2.6\times10^{-2}\) & --- & --- &
--- \\
white sectors N=40 (38) & \(1.8\)--\(2.1\times10^{-2}\) & --- & --- &
--- \\
\end{longtable}
}

Predictions (1) and (2) hold. The residual separates the 4 equilibria
from the non-equilibrium states by 13 orders of magnitude, with not one
intermediate. \(\lambda_{\max}\) matches the prediction exactly at N=5
and runs 4--5\% high at N=40 --- consistent with the N=40 window fit
(\(R^2=0.998\)) including a nonlinear-correction regime. The ratio of
the two N=40 equilibria's \(\lambda_{\max}\), 1.005, matches the
measured burst-rate ratio 1.009. The decisive entry is the last
equilibrium row: \textbf{the inert white-origin parent is a stable
relative equilibrium (zero unstable eigenvalues).} This resolves the
asymmetry of §6 (coherence necessary but not sufficient), and across all
86 tested states (4 equilibria + 42 sectors + 40 generic states),

\[\boxed{\text{latency exponential amplification}\iff\text{unstable relative equilibrium}}\]

holds without exception. The equivalence can be written as a trichotomy:
non-equilibrium states depart immediately by a finite angle, stable
relative equilibria never depart, and only unstable relative equilibria
depart exponentially after a latency. What determines the mode of time
evolution is neither the provenance of the state nor its apparent
coherence, but the dynamical orbit class to which the initial state
belongs.

\subsubsection{5.2 Verification 2: Logarithmic Law of the Perturbation
Amplitude and Measurement of the Intrinsic Floor
(ONS-1)}\label{verification-2-logarithmic-law-of-the-perturbation-amplitude-and-measurement-of-the-intrinsic-floor-ons-1}

A random perturbation \(\varepsilon\eta_\perp\) orthogonal to the parent
plane (\(\eta\) over 3 seeds, \(\varepsilon=10^{-4}\)--\(10^{-14}\)) was
injected onto the control parent. \(t_{\mathrm{cross}}\) is linear in
\(-\ln\varepsilon\) (canonical source:
\texttt{onset\_eps\_sweep\_N*.json}); over \(\varepsilon\ge10^{-12}\)
the slope is 40.35 at N=5 (prediction \(2/\text{rate}_f=40.52\), ratio
0.996, \(R^2=0.9875\)) and 65.45 at N=40 (prediction 57.16, ratio 1.145,
\(R^2=0.9985\)). The N=40 excess is consistent with the initial growth
of a random perturbation being compositionally slower on an equilibrium
carrying 383 unstable modes (growth rates spread over 0.0160--0.0182),
and the measured per-run window rates, which rise monotonically from
0.027 toward 0.035 as \(\varepsilon\) decreases, corroborate this.

At \(\varepsilon=10^{-14}\), \(t_{\mathrm{cross}}\) saturates for all
seeds at the value of the seedless natural trajectory (N=5: 1166, N=40:
2011) --- evidence that the injected perturbation fell below the
system's intrinsic seed. Extrapolating the line, the intrinsic floor is
measured as \(\varepsilon_{\mathrm{eff}}=3.4\times10^{-13}\) (N=5) and
\(2.5\times10^{-13}\) (N=40), of the same order as the parent's
eigenmode residual \(2.14\times10^{-13}\) (Paper 8 §9.1). That is,
\textbf{the ``first minute difference'' that set the latency time of the
seedless trajectory is the eigenmode residual of the parent construction
itself} --- a partial answer, by dynamical measurement, to Paper 8's
open problem (i). Both the pre-registered pooled regression including
the saturation point (\(\varepsilon=10^{-14}\)) and the post-hoc
analysis excluding the saturation regime are recorded side by side in
the canonical source (\texttt{onset\_eps\_sweep\_analysis\_v1.json}).

\subsubsection{5.3 Verification 3: Direct Injection of Generic
Zero-Closure States
(ONS-3)}\label{verification-3-direct-injection-of-generic-zero-closure-states-ons-3}

The white sectors of §4 are 42 generator-derived states, and since the
dynamics is nonlinear, the behavior of other points on the closure
manifold does not follow from their results. Two ensembles were
therefore injected directly (10 states each per N; canonical source:
\texttt{onset\_general\_closed\_N*.json}).

\begin{itemize}
\tightlist
\item
  \textbf{A: direct random samples from the manifold} ---
  \(Z=(u_1+iu_2)/\sqrt2\) from the two QR-orthonormal columns of
  \(M\times2\) Gaussian white noise. Closure holds identically (measured
  \(\le5.6\times10^{-16}\)).
\item
  \textbf{B: random mixtures of admissible sectors} --- the 42 sectors
  superposed with complex Gaussian coefficients, then projected onto the
  closure manifold by minimal change (SVD singular-value equalization of
  \([X\,Y]\); projection distances 0.003--0.16 recorded; post-projection
  closure \(\le3.9\times10^{-16}\)).
\end{itemize}

Result: \textbf{all 40 states are of the immediate type} (crossing
11--21), with no latency-burst and no ``other.'' As predicted. The
dichotomy is not a peculiarity of the single-harmonic-sector
construction; it is a consequence of the state structure --- equilibrium
or not.

\subsection{6. Provenance Independence}\label{provenance-independence}

The white-origin parent results show that the boundary of the dichotomy
is not ``random origin vs.~deterministic construction.'' Starting from
the same white noise, if the self-consistency iteration converges to an
eigenmode (N=40), that state shows the same latency burst at the same
rate as the deterministic control. At N=5, meanwhile, the iteration
converged to an eigen-solution of a different family and did not amplify
out to 12000 steps. That is:

\begin{itemize}
\tightlist
\item
  The first boundary determining the onset mode is whether the state is
  a \textbf{self-consistent relative equilibrium}, and what finally
  determines the presence of amplification is \textbf{the stability of
  that relative equilibrium} (§5.1). How the state was prepared plays no
  role at either boundary.
\item
  Among self-consistent closures there are amplifying solutions
  (unstable equilibria) and inert solutions (stable equilibria).
  Coherence alone is a \textbf{necessary but not sufficient condition}
  for amplification.
\end{itemize}

This asymmetry does not weaken the claim of this paper, whose
implication runs in a fixed direction: if amplification was observed,
the initial state was a self-consistent closure. That there is not one
counterexample among the 82 non-equilibrium states (42 sectors + 40
generic states, §5.3) is the measured basis of this implication.
Moreover, the tangent-spectrum measurement of §5.1 resolved the
asymmetry itself: the inert white-origin parent (N=5) is a stable
relative equilibrium (zero unstable eigenvalues), and what finally
separates amplification from its absence is whether the state is an
equilibrium and, if so, its stability.

\subsection{7. Record of Refutation and
Correction}\label{record-of-refutation-and-correction}

This result was not obtained in a single experiment. The full process is
recorded.

\begin{enumerate}
\def\labelenumi{\arabic{enumi}.}
\tightlist
\item
  \textbf{Defect of the generator's first version (v2).} Writing the row
  closure in DFT form, \(\sum_n w^2=N\sum_k c_kc_{N-k}\): v2 satisfied
  closure through the \textbf{cancellation of cross pairs} of harmonics
  \(k\) and \(N-k\), leaving the self-paired sectors (DC, Nyquist)
  unclosed. Based on the theorem of §2.2, v3 projects them out at the
  source of the noise. v3 consumes random numbers in the same order as
  v2; the parent vector and raw noise are bitwise identical.
\item
  \textbf{Measurement error on the experimental side.} The first E-M9r
  run used a wrong quadratic form for the sectors' closure measurement
  and misreported ``white sectors are unclosed.'' Under the correct
  field closure, the 42 white sectors are legitimate, exactly
  zero-closed states (conservation verified: drift \(\sim10^{-16}\) over
  500 steps).
\item
  \textbf{Correction via control experiments.} The generator was copied
  into a managed area with SHA-256 identity, and the erroneous results
  were bit-reproduced with identical seeds before being corrected. We do
  not claim a correction without reproducing the error.
\item
  \textbf{Survival of the dichotomy.} The central dichotomy (latency
  burst vs.~immediate) was observed in both v2 (defective) and v3
  (corrected). The bookkeeping error changed the handling of the
  self-paired sectors but did not change the dichotomy itself.
\end{enumerate}

\subsection{8. Boundary Conditions (What This Paper Does Not
Prove)}\label{boundary-conditions-what-this-paper-does-not-prove}

\begin{enumerate}
\def\labelenumi{\arabic{enumi}.}
\tightlist
\item
  \textbf{The theoretical classification of stability is incomplete.}
  That amplification is decided by the stability of the equilibrium was
  measured in §5.1 (inert parent = stable equilibrium,
  \(\lambda_{\max}=0\)). But no theoretical classification is given that
  predicts, from the structure of a solution (N, family, spectrum),
  which self-consistent families are stable and which unstable. Why the
  N=5 white-origin solution converged to a stable one and the N=40
  solution to an unstable one remains open.
\item
  \textbf{Seed count.} One generator seed per \(N\) (N=5: 2, N=40: 1).
  On the non-equilibrium side, however, there was no exception among the
  42 independent sectors within the same seed plus the 40 generic
  zero-closure states from independent random streams (§5.3, rng
  94000+/95000+). A sweep over generator seeds has not been performed.
\item
  \textbf{Time range.} The profile measurement uses \(T=6000\); the v3
  main run \(T=12000\). The ``absence of amplification'' of the white
  sectors is definitive as a finite departure within the first 2 steps,
  but a second amplification beyond the observation window is not
  logically excluded.
\item
  \textbf{Range of \(N\).} Only N=5 and 40. The N=300 of Paper 8's main
  series was not run in this control.
\item
  \textbf{The capture problem.} The spontaneous formation of a
  self-consistent closure out of the white sea is not observed in the
  linear dynamics of this paper. This is the province of the inelastic
  generation channels from Paper 9 {[}5{]} onward.
\item
  \textbf{Limitation of terminology.} ``Inflation'' and
  ``amplification'' in this paper are phenomenon names of this series
  (the time structure latency \(\to\) geometric \(\to\) metastable of
  \(f\)), not cosmological inflation (exponential expansion of the
  metric) itself. The correspondence of §10 is a structural analogy, not
  a dynamical derivation in field theory or general relativity.
\item
  \textbf{The linear no-generation hypothesis is unproven.} Of the
  mechanistic reading of §5, the absence of an explicit generation
  vertex is a verified fact, and relative equilibrium, instability, the
  logarithmic law, and the intrinsic floor became measurements in
  §5.1--5.3. But that iteration of the state-dependent unitary cannot
  create new harmonic content (preservation of the invariant subspace
  corresponding to the initial spectral support) remains unproven. The
  central results --- the dichotomy of §4 and the equivalence of §5.1
  --- are measurements and do not depend on this unproven part.
\end{enumerate}

\subsection{9. Conclusion}\label{conclusion}

We executed the systematic control left by Paper 8 --- starting from
generic states satisfying only the closure condition. The result is a
complete dichotomy of onset modes. Only series starting from
self-consistent closures (single coherent condensates) show latency and
18 decades of geometric amplification. The 82 non-equilibrium
zero-closure states (42 white sectors + 40 generic states) have, without
exception, no amplification epoch and move immediately, by finite-angle
rotation, to the metastable three-directional structure. The presence of
amplification does not depend on the provenance of the state; it is
decided by the structure of the state alone.

The mechanism closed with measurements. The one-step residual separates
equilibria from non-equilibria by 13 orders of magnitude; the maximal
tangent-map growth rate equals half the measured burst rate (exactly at
N=5, within 4\% at N=40); the inert parent is a stable equilibrium (zero
unstable eigenvalues); the logarithmic law of the perturbation amplitude
matches the prediction down to its slope; and its saturation measures
the intrinsic floor of the seedless trajectory at the location of the
eigenmode residual. Across all 86 tested states, the equivalence ---
latency exponential amplification \(\iff\) unstable relative equilibrium
--- holds without exception.

Together with Papers 7 and 8, the causal ledger closes in the following
form. The explicit seeds are not the cause of the phenomenon (Paper 8).
The cause is that the initial state belongs to an unstable
self-consistent relative equilibrium (this paper). The necessary and
sufficient condition for latency exponential amplification is, within
the tested range, an unstable self-consistent relative equilibrium. The
metastable three-directional structure is a terminal point independent
of the onset mode, and only the trace of geometric amplification remains
as the fossil showing that the system began from an unstable relative
equilibrium --- a coherent condensate.

What remains unresolved: (i) the theoretical classification of the
correspondence between self-consistent solution families and stability
(why the N=5 white-origin solution is stable and the N=40 solution
unstable); (ii) the capture from white chaos to condensate; (iii) the
inelastic generation of structure (matter) from the condensate --- (iii)
is already underway in the series deriving the vertex from closure
conservation {[}5{]}.

\subsection{10. Relation to Prior Work (via Re-Survey After the
Design)}\label{relation-to-prior-work-via-re-survey-after-the-design}

The control experiments of this paper were not designed as replication
or refutation of prior work. The starting point was a question internal
to this series --- ``would not a flat spectrum with all harmonics mixed
be a more natural initial state than concentration on a single
wavelength at the minimal resolution?'' --- and the working hypothesis
at design time was ``the white sea should show the same geometric
evolution.'' The result is the rejection of that working hypothesis
(§4). Having obtained a result contrary to expectation, we re-surveyed
the literature on initial conditions in inflationary cosmology, and
found the following.

\textbf{The standard theory does not predict the initial state; it
assumes it.} The starting point of slow-roll inflation {[}6,7{]} is a
nearly homogeneous, potential-dominated inflaton condensate with
Bunch--Davies vacuum fluctuations on top. The origin of this initial
state is undecided --- the ``initial conditions problem of inflation''
{[}8,9{]} --- and the observed nearly flat spectrum (\(n_s=0.9649\)
{[}14{]}) is an \textbf{output} of inflation, not an initial condition.
That is, the standard theory also posits, without deriving it, the onset
mode of a coherent condensate. The measurements of this paper show,
dynamically and inside this model, that only the state corresponding to
that assumption (the self-consistent closure) amplifies.

\textbf{The claim that ``inflation ignites even from inhomogeneous
initial conditions'' is a model-dependent numerical result of a few
groups.} Numerical-relativity studies {[}10,11{]} reported that in
large-field (plateau-type) models inflation begins even from
gradient-energy-dominated inhomogeneous initial conditions. But the
initial data of these computations restrict the field values to the
potential plateau before adding inhomogeneity --- in the language of
this paper, they correspond not to a ``white sea'' but to a ``condensate
plus perturbations.'' The coherence assumption is not removed; it is
relocated into the construction of the initial data. If this reading is
correct, those results do not contradict the dichotomy of this paper (no
amplification from generic states).

\textbf{A decision criterion for the critics' intuition.} Against the
criticism that inflation demands initial conditions more special than
the problem it is supposed to solve {[}12,13{]}, and the resulting
standoff (called a ``schism'' by the participants themselves {[}13{]}),
the results of this paper provide a dynamical criterion of the following
form. First, as a result inside the model: trajectories with the
exponential-amplification time structure did not arise directly from
white non-equilibrium states, but only from unstable self-consistent
relative equilibria (§4--§5, no exception among 86 states). Translated,
as an analogy, into cosmological language: an onset state bearing the
trace of exponential amplification cannot be white chaos and must be an
unstable coherent condensate --- with the proviso that the range of this
translation extends only to the structural analogy of this model, and
any transfer into field theory is subject to the limitation of §8.6.

The order of emphasis is this. This paper did not draw its hypotheses
from these references. An independently designed control returned an
unexpected dichotomy, and the post-hoc re-survey revealed that this
dichotomy amounts to a decision experiment, within this model, on a
point of controversy that has remained open for forty years --- is
coherence necessary for the onset? Although the model contains no field
theory, it arrived at a question of the same structure as the
initial-conditions problem and returned a definite answer there
(coherence is necessary). This convergence is itself a verifiable
outcome of the method of this series: fix the axioms minimally and let
the dynamics answer.

\subsection{11. Reproducibility}\label{reproducibility}

\begin{itemize}
\tightlist
\item
  The dynamics code is the Paper 7 original
  (\texttt{run\_preliminary\_seed\_ablation\_v1.py}), imported read-only
  (SHA-256: \texttt{75a10a5b951302be…}). Generator v3 SHA-256:
  \texttt{d3217579ab54ce29…}. Both hashes are recorded in the
  \texttt{imports} of every result JSON.
\item
  Canonical sources for the generator, classifier, and audits:
  \texttt{make\_parent\_white\_managed\_v1/} (v2 control copy, v3, v4,
  unit tests, classifier control). v2 and v3 are bitwise identical in
  parent vector and raw noise at identical seeds. The control
  experiments include SHA-256 identity with the Codex original and bit
  reproduction of the erroneous results.
\item
  Canonical sources for the experiments:
  \texttt{paper8\_em9r\_white\_harmonics\_inflation\_v1/} --- all JSONs
  and figures of the first run (v1), the bookkeeping correction
  (fix\_v2), the v3-generator run (v3\_result), and the profile
  measurement (profile; canonical source of Table 1 and Figures 1--2).
\item
  Canonical sources for the mechanism verification:
  \texttt{paper8\_onset\_mechanism\_v1/} --- ONS-1 (\(\varepsilon\)
  sweep: \texttt{onset\_eps\_sweep\_N*.json}, post-hoc analysis
  \texttt{onset\_eps\_sweep\_analysis\_v1.json}), ONS-2 (residual and
  tangent spectrum: \texttt{onset\_equilibrium\_residual\_N*.json},
  canonical source of Table 2), ONS-3 (generic-state injection:
  \texttt{onset\_general\_closed\_N*.json}) --- all scripts, JSONs, and
  figures. Criteria and predictions are fixed before execution at the
  head of each script. Random seeds: \(\eta\)=93000+j, direct
  samples=94000+j, mixtures=95000+j, injection auxiliary=96000+series.
\item
  Commit series: \texttt{52c83f13} (initial-state classification)
  \(\to\) \texttt{aeaf8e57} (first E-M9r) \(\to\) \texttt{d9346e37}
  (bookkeeping correction) \(\to\) \texttt{e5232885} (v3 re-run) \(\to\)
  \texttt{94aa76f9} (profile measurement) \(\to\) \texttt{b1cded81} (v4
  generator, catalog closure audit) \(\to\) \texttt{e6681744}
  (discussion note) \(\to\) \texttt{2f443467} (mechanism verification
  ONS-1/2/3).
\item
  The classification criteria (§3) are stated as pre-registered at the
  head of the measurement programs and saved in the \texttt{criteria} of
  the result JSONs. All random seeds are recorded (generator: N=5
  \(\to\) 2, N=40 \(\to\) 1; injection auxiliary: parent 91000, sector
  \(92000+k\)).
\item
  Environment: Python 3.9.6, NumPy 2.0.2, macOS arm64.
\end{itemize}

\subsection{References}\label{references}

{[}1{]} N. Kihara, Foundational Axiom System of the Anonymous
Equal-Amplitude Composite Wave Model (Purified Definition Paper),
Zenodo. Concept DOI:
\href{https://doi.org/10.5281/zenodo.21315735}{10.5281/zenodo.21315735}.

{[}2{]} N. Kihara, Arrest of Spontaneous Splitting and Creation of
Additional Axes in Closed N-Body Relational Wave Systems (Paper 7),
Zenodo. Concept DOI:
\href{https://doi.org/10.5281/zenodo.21543070}{10.5281/zenodo.21543070}.

{[}3{]} N. Kihara, Causal Separation of the Time Structure of
Three-Direction Formation in Closed N-Body Relational Wave Systems via
Two-Stage Seed Removal (Paper 8), Zenodo. Concept DOI:
\href{https://doi.org/10.5281/zenodo.21614402}{10.5281/zenodo.21614402}.

{[}4{]} N. Kihara, The Geometric Identity of the Zero-Square-Sum
Constraint and Scale Invariance (Explanatory Note), Zenodo. Concept DOI:
\href{https://doi.org/10.5281/zenodo.21495305}{10.5281/zenodo.21495305}.

{[}5{]} N. Kihara, The Generation Structure of Fermions (Paper 9),
Zenodo. Concept DOI:
\href{https://doi.org/10.5281/zenodo.21766706}{10.5281/zenodo.21766706}.

{[}6{]} A. H. Guth, Inflationary universe: A possible solution to the
horizon and flatness problems, Phys. Rev.~D \textbf{23}, 347 (1981).

{[}7{]} A. D. Linde, Chaotic inflation, Phys. Lett. B \textbf{129}, 177
(1983).

{[}8{]} R. Brandenberger, Initial conditions for inflation --- A short
review, Int. J. Mod. Phys. D \textbf{26}, 1740002 (2017).
\href{https://arxiv.org/abs/1601.01918}{arXiv:1601.01918}.

{[}9{]} A. Linde, On the problem of initial conditions for inflation,
Found. Phys. \textbf{48}, 1246 (2018).
\href{https://arxiv.org/abs/1710.04278}{arXiv:1710.04278}.

{[}10{]} W. E. East, M. Kleban, A. Linde, L. Senatore, Beginning
inflation in an inhomogeneous universe, JCAP \textbf{09} (2016) 010.
\href{https://arxiv.org/abs/1511.05143}{arXiv:1511.05143}.

{[}11{]} K. Clough, E. A. Lim, B. S. DiNunno, W. Fischler, R. Flauger,
S. Paban, Robustness of inflation to inhomogeneous initial conditions,
JCAP \textbf{09} (2017) 025.
\href{https://arxiv.org/abs/1608.04408}{arXiv:1608.04408}.

{[}12{]} A. Ijjas, P. J. Steinhardt, A. Loeb, Inflationary paradigm in
trouble after Planck2013, Phys. Lett. B \textbf{723}, 261 (2013).
\href{https://arxiv.org/abs/1304.2785}{arXiv:1304.2785}.

{[}13{]} A. Ijjas, P. J. Steinhardt, A. Loeb, Inflationary schism, Phys.
Lett. B \textbf{736}, 142 (2014).
\href{https://arxiv.org/abs/1402.6980}{arXiv:1402.6980}.

{[}14{]} Planck Collaboration, Planck 2018 results. X. Constraints on
inflation, Astron. Astrophys. \textbf{641}, A10 (2020).
\href{https://arxiv.org/abs/1807.06211}{arXiv:1807.06211}.

\end{document}
